% Part: proof-theory
% Chapter: natural-deduction
% Section: quantifiers

\documentclass[../../../include/open-logic-section]{subfiles}

\begin{document}

\olfileid{pt}{nat}{qua}

\olsection{\usetoken{P}{proof} منظم و جانشینی}

قواعد سورها به سبب «شرط‌های متغیر ویژه» که بر قواعد
\Elim{\lexists} و~\Intro{\lforall} اعمال می‌شوند پیچیدگی‌هایی پدید
می‌آورند. بیشتر این پیچیدگی‌ها را می‌توان با تضمین این امر برطرف کرد
که !!{constant}~$c$ در این استنتاج‌ها هرگز دوباره به کار نرود. تعریف
زیر را معرفی کنیم.

\begin{defn}
یک !!{proof} استنتاج طبیعی~$\delta$ \emph{منظم} است اگر:
\begin{enumerate}
  \item هیچ متغیر ویژهٔ~$a$ به‌عنوان متغیر ویژهٔ بیش از یک استنتاج
  \Elim{\lexists} یا~\Intro{\lforall} به کار نرود؛ و
  \item اگر $a$ متغیر ویژهٔ یک استنتاج \Elim{\lexists} یا
  ~\Intro{\lforall} باشد، به‌ترتیب فقط در زیر-!!{proof} بی‌واسطه‌ای
  که به مقدمهٔ فرعی می‌رسد، یا در زیر-برهان بی‌واسطه‌ای که به مقدمه
  می‌رسد، ظاهر شود.
\end{enumerate}
\end{defn}

\begin{lem}\ollabel{lem:subst-regular}
اگر $\delta$ منظم باشد و به~$!C$ پایان یابد، $c$ یک !!a{constant}
باشد که در~$\delta$ به‌عنوان متغیر ویژه به کار نرفته است، و~$t$ ترمی
باشد که هیچ متغیر ویژهٔ~$\delta$ را در بر ندارد، آنگاه
$\Subst{\delta}{t}{c}$ که با جایگزینی هر رخداد~$c$ در~$\delta$ با~$t$
به دست می‌آید، یک !!{proof} منظم است. !!{formula} پایانی
$\Subst{\delta}{t}{c}$ برابر $\Subst{!C}{t}{c}$ است. اگر $!A^x$ فرض
باز~$\delta$ باشد، آنگاه $\Subst{!A}{t}{c}^x$ فرض باز
$\Subst{\delta}{t}{c}$ است.
\end{lem}

\begin{proof}
با استقرا بر $\pheight{\delta}$. اگر $\pheight{\delta}=0$، برهان فقط
از فرض~$!A^x$ تشکیل شده است. آنگاه
$\Subst{\delta}{t}{c}$ همان~$\Subst{!A}{t}{c}^x$ است.

اگر $\pheight{\delta}>0$، برحسب آخرین استنتاج~$\delta$ حالت‌ها را جدا
می‌کنیم. حالت‌های مربوط به قواعد رابط‌های گزاره‌ای آسان‌اند و به‌عنوان
تمرین واگذار می‌شوند. حالت‌هایی را بررسی می‌کنیم که $\delta$ به یک
استنتاج سوردار پایان می‌یابد.
  \begin{enumerate}
    \item اگر آخرین استنتاج $\Intro{\lforall}$ باشد، $\delta$ شکل زیر
    را دارد:
    \[
    \AxiomC{}
    \RightLabel{$\delta'$}
    \DeduceC{$!A(a,c)$}
    \RightLabel{\Intro{\lforall}}
    \UnaryInfC{$\lforall[x][!A(x,c)]$}
    \DisplayProof
    \]
    بنا بر فرض استقرا، $\Subst{\delta'}{t}{c}$ یک !!{proof} منظم برای
    $!A(a,t)$ است. چون $a$ متغیر ویژهٔ~$\delta$ است، در~$t$ ظاهر
    نمی‌شود. بنابراین آخرین استنتاج در $\Subst{\delta}{t}{c}$،
    \[
    \AxiomC{}
    \RightLabel{$\Subst{\delta'}{t}{c}$}
    \DeduceC{$!A(a,t)$}
    \RightLabel{\Intro{\lforall}}
    \UnaryInfC{$\lforall[x][!A(x,t)]$}
    \DisplayProof
    \]
    درست است: چون $t$ شامل~$a$ نیست، نه نتیجه و نه هیچ فرض بازی شامل
    ~‌$a$ نیست.

    \item اگر آخرین استنتاج $\Elim{\lforall}$ باشد، $\delta$ شکل زیر
    را دارد:
    \[
    \AxiomC{}
    \RightLabel{$\delta'$}
    \DeduceC{$\lforall[x][!A(x,c)]$}
    \RightLabel{\Elim{\lforall}}
    \UnaryInfC{$!A(s(c),c)$}
    \DisplayProof
    \]
    بنا بر فرض استقرا، $\Subst{\delta'}{t}{c}$ یک !!{proof} منظم برای
    $\lforall[x][!A(x,t)]$ است. (ترم~$s(c)$ در نتیجه ممکن است شامل~$c$
    باشد یا نباشد.) آخرین استنتاج در $\Subst{\delta}{t}{c}$،
    \[
    \AxiomC{}
    \RightLabel{$\Subst{\delta'}{t}{c}$}
    \DeduceC{$\lforall[x][!A(x,t)]$}
    \RightLabel{\Elim{\lforall}}
    \UnaryInfC{$!A(s(t),t)$}
    \DisplayProof
    \]
    درست است و
    $!A(s(t),t)\ident\Subst{!A(s(c),c)}{t}{c}$.
    حالت \Intro{\lexists} مشابه است.

  \item اگر آخرین استنتاج $\Elim{\lexists}$ باشد، $\delta$ شکل زیر
    را دارد:
    \[
    \AxiomC{}
    \RightLabel{$\delta_1'$}
    \DeduceC{$\lexists[x][!A(x,c)]$}
    \AxiomC{$\Discharge{!A(a,c)}{x}$}
    \RightLabel{$\delta_2'$}
    \DeduceC{$!C$}
    \DischargeRule{\Elim{\lexists}}{x}
    \BinaryInfC{$!C$}
    \DisplayProof
    \]
    بنا بر فرض استقرا، $\Subst{\delta_1'}{t}{c}$ و
    $\Subst{\delta_2'}{t}{c}$ !!{proof}های منظمی هستند، به‌ترتیب برای
    $\lexists[x][!A(x,t)]$ و برای $\Subst{!C}{t}{c}$ از فرض باز
    ~‌$!A(a,t)^x$. چون $a$ متغیر ویژهٔ~$\delta$ است، در~$t$ ظاهر
    نمی‌شود. بنابراین آخرین استنتاج در $\Subst{\delta}{t}{c}$ چنین
    است:
    \[
    \AxiomC{}
    \RightLabel{$\Subst{\delta_1'}{t}{c}$}
    \DeduceC{$\lexists[x][!A(x,t)]$}
    \AxiomC{$\Discharge{!A(a,t)}{x}$}
    \RightLabel{$\Subst{\delta_2'}{t}{c}$}
    \DeduceC{$\Subst{!C}{t}{c}$}
    \DischargeRule{\Elim{\lexists}}{x}
    \BinaryInfC{$\Subst{!C}{t}{c}$}
    \DisplayProof
    \]
    و درست است: داریم
    $\Subst{!A(x,t)}{a}{x}\ident\Subst{!A(a,c)}{t}{c}$، و $a$ نه در
    نتیجهٔ $\Subst{!C}{t}{c}$ و نه در هیچ فرض باز، جز فرض
    مرخص‌شوندهٔ~$!A(a,t)^x$، ظاهر می‌شود.
  \end{enumerate}
\end{proof}

\begin{prop}\ollabel{prop:regular}
اگر $\delta$ !!a{proof}ی در \Log{N1c} یا \Log{N1i} باشد، آنگاه یک
!!{proof} منظم~$\delta'$ با همان ساختار، و به‌ویژه همان !!{height}،
وجود دارد که فقط با جایگزینی متغیرهای ویژه با~$\delta$ تفاوت دارد.
\end{prop}

\begin{proof}
فقط برای مقاصد این اثبات، یک استنتاج متغیر ویژهٔ~$\delta$ را
\emph{پاک} بنامید اگر متغیر ویژهٔ آن، متغیر ویژهٔ استنتاج دیگری نباشد
و جز در زیر-!!{proof} بی‌واسطهٔ همان استنتاج در هیچ جای~$\delta$ ظاهر
نشود؛ در غیر این صورت آن را \emph{ناپاک} بنامید. بگذارید $n$ شمار
استنتاج‌های متغیر ویژهٔ ناپاک در~$\delta$ باشد. با استقرا بر~$n$ نشان
می‌دهیم $\delta$ را می‌توان به !!{proof} منظم مطلوب~$\delta'$ تبدیل
کرد. اگر $n=0$، همهٔ استنتاج‌های متغیر ویژه در~$\delta$ پاک‌اند؛ پس
$\delta$ منظم است و چیزی برای اثبات باقی نمی‌ماند. در غیر این صورت،
یکی از بالاترین استنتاج‌های متغیر ویژه، مثلاً \Intro{\lforall}، را
برگزینید. زیر-!!{proof}~$\delta_1$ از~$\delta$ که به آن پایان می‌یابد
شکل زیر را دارد:
    \[
    \AxiomC{}
    \RightLabel{$\delta_2$}
    \DeduceC{$!A(c)$}
    \RightLabel{\Intro{\lforall}}
    \UnaryInfC{$\lforall[x][!A(x)]$}
    \DisplayProof
    \]
چون $c$ متغیر ویژه است، در~$\lforall[x][!A(x)]$ یا در هیچ فرض بازی
ظاهر نمی‌شود. برهان~$\delta_2$ منظم است، زیرا اصلاً هیچ استنتاج متغیر
ویژه‌ای ندارد. بگذارید $a$ یک !!a{constant} باشد که در~$\delta$ ظاهر
نمی‌شود. بنا بر \olref{lem:subst-regular}،
$\Subst{\delta_2}{a}{c}$ یک برهان منظم برای $!A(a)$ است. بنا بر شرط
متغیر ویژه، هیچ فرض باز~$\delta_2$ شامل~$c$ نیست؛ پس هیچ فرض باز
$\Subst{\delta_2}{a}{c}$ شامل~$a$ نیست. بنابراین می‌توانیم
~\Intro{\lforall} را اعمال کنیم. اگر $\delta_1$ را در~$\delta$ با
برهان~$\delta_1'$ زیر جایگزین کنیم،
    \[
    \AxiomC{}
    \RightLabel{$\Subst{\delta_2}{a}{c}$}
    \DeduceC{$!A(a)$}
    \RightLabel{\Intro{\lforall}}
    \UnaryInfC{$\lforall[x][!A(x)]$}
    \DisplayProof
    \]
یک استنتاج متغیر ویژهٔ ناپاک را با استنتاجی پاک جایگزین کرده‌ایم، و
تنها تفاوت، جایگزینی متغیر ویژهٔ~$c$ با~$a$ است. برهان حاصل $n-1$
استنتاج متغیر ویژهٔ ناپاک دارد و نتیجه از فرض استقرا به دست می‌آید.
\end{proof}

\begin{prob}
حالتی از اثبات \olref{prop:regular} را شرح دهید که بالاترین استنتاج
متغیر ویژه~\Elim{\lexists} است.
\end{prob}

به گزاره‌ای که همین الآن اثبات شد صریحاً استناد نخواهیم کرد، اما فرض
می‌کنیم همهٔ !!{proof}های مورد نظر منظم‌اند؛ اگر نباشند، همیشه می‌توان
با جایگزینی متغیرهای ویژه آن‌ها را منظم کرد.

\Olref{lem:subst-regular} و \olref{prop:regular} برای
\Log{N2c} و~\Log{N2i} نیز برقرارند. اثبات‌ها را به‌عنوان تمرین
واگذار می‌کنیم.

\end{document}
